% GENERATED FILE. Edit canonical lessons, then run npm run generate:books.
% canonical-source-hash: fnv1a64:0cca52434b187a2a
% canonical-lessons: MW-C147-mehanti, MW-C147-imandar, MW-C147-dayalu, MW-C147-kanjus, MW-C147-bhalo

\chapter{Hard-working, Honest, Kind, Miserly, Good-natured}
\label{ch:mw-hard-working-honest-kind-miserly-good-natured}
\hlchaptermodality{147}

\begin{chapteropening}
\textbf{By the end of this chapter:} \emph{I can say hard-working, honest, kind, miserly and good-natured.}

Complete the last lesson of chapter 147: I can say hard-working, honest, kind, miserly and good-natured.
\end{chapteropening}

\section[\texorpdfstring{\mw{मेहनती} (mehantī)}{मेहनती (mehantī)}]{\mw{मेहनती} (mehantī) — hard-working}

\label{lesson:MW-C147-mehanti}



\begin{quote}
Before the new one: say the Marwadi for silly, then the Marwadi for lazy.
\end{quote}

\subsection*{You'll want to know: \mw{मेहनती}}
\textbf{\mw{मेहनती}} — \emph{mehantī} — \textquotedblleft{}hard-working\textquotedblright{}.

\begin{cousinweb}
\textbf{\mw{मेहनती}} (\emph{mehantī}) is a loan from Arabic \emph{miḥna}, \textquotedblleft{}trial, ordeal\textquotedblright{}, which reached Marwadi through Persian. Its Arabic root is \emph{m-ḥ-n} (Lane). Persian \emph{mihnat} means \textquotedblleft{}toil\textquotedblright{}; Marwadi added the ending \emph{-ī}.
\end{cousinweb}

\begin{grammarlens}[title={What you've built}]
One more word for this chapter.
\end{grammarlens}

\subsection*{Your turn}
\begin{itemize}
\raggedright
  \item \emph{Say it:} \emph{mehantī}
  \item \emph{Say it:} \emph{mehantī}, once more
  \item \emph{Say it:} say \emph{mehantī}
  \item \emph{Recall:} say the Marwadi for silly, then the Marwadi for lazy, then say \emph{mehantī} again
\end{itemize}

\begin{tcolorbox}[breakable,colback=teal!4,colframe=teal!35!black,title={Before you move on}]
What does \mw{मेहनती} mean? (\textquotedblleft{}Hard-working\textquotedblright{}.) Say it once more.
\end{tcolorbox}
\section[\texorpdfstring{\mw{ईमानदार} (īmāndār)}{ईमानदार (īmāndār)}]{\mw{ईमानदार} (īmāndār) — honest}

\label{lesson:MW-C147-imandar}



\begin{quote}
Before the new one: say the Marwadi for lazy, then the Marwadi for hard-working.
\end{quote}

\subsection*{You'll want to know: \mw{ईमानदार}}
\textbf{\mw{ईमानदार}} — \emph{īmāndār} — \textquotedblleft{}honest\textquotedblright{}.

\begin{cousinweb}
\textbf{\mw{ईमानदार}} (\emph{īmāndār}) joins Arabic \emph{īmān}, \textquotedblleft{}faith\textquotedblright{}, whose root is \emph{ʾ-m-n} (Lane), with Persian \emph{-dār}, \textquotedblleft{}holder\textquotedblright{}: \textquotedblleft{}one who keeps faith\textquotedblright{}. It reached Marwadi through Persian.
\end{cousinweb}

\begin{grammarlens}[title={What you've built}]
One more word for this chapter.
\end{grammarlens}

\subsection*{Your turn}
\begin{itemize}
\raggedright
  \item \emph{Say it:} \emph{īmāndār}
  \item \emph{Say it:} \emph{īmāndār}, once more
  \item \emph{Say it:} say \emph{īmāndār}
  \item \emph{Recall:} say the Marwadi for lazy, then the Marwadi for hard-working, then say \emph{īmāndār} again
\end{itemize}

\begin{tcolorbox}[breakable,colback=teal!4,colframe=teal!35!black,title={Before you move on}]
What does \mw{ईमानदार} mean? (\textquotedblleft{}Honest\textquotedblright{}.) Say it once more.
\end{tcolorbox}
\section[\texorpdfstring{\mw{दयाळु} (dayāḷu)}{दयाळु (dayāḷu)}]{\mw{दयाळु} (dayāḷu) — kind}

\label{lesson:MW-C147-dayalu}



\begin{quote}
Before the new one: say the Marwadi for hard-working, then the Marwadi for honest.
\end{quote}

\subsection*{You'll want to know: \mw{दयाळु}}
\textbf{\mw{दयाळु}} — \emph{dayāḷu} — \textquotedblleft{}kind\textquotedblright{}.

\begin{cousinweb}
\textbf{\mw{दयाळु}} (\emph{dayāḷu}) is a semi-tatsama: Sanskrit \emph{dayālu}, \textquotedblleft{}compassionate\textquotedblright{}, borrowed after the old sound changes had run their course and then lightly reshaped. Marwadi says the \emph{l} curled back; CDIAL has no entry for the word, only for its base \emph{dayā}, \textquotedblleft{}sympathy\textquotedblright{} (CDIAL 6185).
\end{cousinweb}

\begin{grammarlens}[title={What you've built}]
One more word for this chapter.
\end{grammarlens}

\subsection*{Your turn}
\begin{itemize}
\raggedright
  \item \emph{Say it:} \emph{dayāḷu}
  \item \emph{Say it:} \emph{dayāḷu}, once more
  \item \emph{Say it:} say \emph{dayāḷu}
  \item \emph{Recall:} say the Marwadi for hard-working, then the Marwadi for honest, then say \emph{dayāḷu} again
\end{itemize}

\begin{tcolorbox}[breakable,colback=teal!4,colframe=teal!35!black,title={Before you move on}]
What does \mw{दयाळु} mean? (\textquotedblleft{}Kind\textquotedblright{}.) Say it once more.
\end{tcolorbox}
\section[\texorpdfstring{\mw{कंजूस} (kañjūs)}{कंजूस (kañjūs)}]{\mw{कंजूस} (kañjūs) — miserly}

\label{lesson:MW-C147-kanjus}



\begin{quote}
Before the new one: say the Marwadi for honest, then the Marwadi for kind.
\end{quote}

\subsection*{You'll want to know: \mw{कंजूस}}
\textbf{\mw{कंजूस}} — \emph{kañjūs} — \textquotedblleft{}miserly\textquotedblright{}.

\begin{cousinweb}
The origin of \textbf{\mw{कंजूस}} (\emph{kañjūs}) is not settled. CDIAL, searched by form and by meaning, has no entry for it; it is often said to come from Persian.
\end{cousinweb}

\begin{grammarlens}[title={What you've built}]
One more word for this chapter.
\end{grammarlens}

\subsection*{Your turn}
\begin{itemize}
\raggedright
  \item \emph{Say it:} \emph{kañjūs}
  \item \emph{Say it:} \emph{kañjūs}, once more
  \item \emph{Say it:} say \emph{kañjūs}
  \item \emph{Recall:} say the Marwadi for honest, then the Marwadi for kind, then say \emph{kañjūs} again
\end{itemize}

\begin{tcolorbox}[breakable,colback=teal!4,colframe=teal!35!black,title={Before you move on}]
What does \mw{कंजूस} mean? (\textquotedblleft{}Miserly\textquotedblright{}.) Say it once more.
\end{tcolorbox}
\section[\texorpdfstring{\mw{भलो} (bhalo)}{भलो (bhalo)}]{\mw{भलो} (bhalo) — good-natured}

\label{lesson:MW-C147-bhalo}



\begin{quote}
Before the new one: say the Marwadi for kind, then the Marwadi for miserly.
\end{quote}

\subsection*{You'll want to know: \mw{भलो}}
\textbf{\mw{भलो}} — \emph{bhalo} — \textquotedblleft{}good-natured\textquotedblright{}.

\begin{cousinweb}
\textbf{\mw{भलो}} (\emph{bhalo}) is a tadbhava, a word inherited from Sanskrit and reshaped by sound change. It comes from \emph{bhalla}, \textquotedblleft{}auspicious\textquotedblright{}. CDIAL lists Old Marwari \emph{bhalaü} under it (CDIAL 9408). CDIAL builds the Sanskrit word on \emph{bhadra}, \textquotedblleft{}good\textquotedblright{}.
\end{cousinweb}

\begin{grammarlens}[title={What you've built}]
One more word for this chapter.
\end{grammarlens}

\subsection*{Your turn}
\begin{itemize}
\raggedright
  \item \emph{Say it:} \emph{bhalo}
  \item \emph{Say it:} \emph{bhalo}, once more
  \item \emph{Say it:} say \emph{bhalo}
  \item \emph{Recall:} say the Marwadi for kind, then the Marwadi for miserly, then say \emph{bhalo} again
\end{itemize}

\begin{tcolorbox}[breakable,colback=teal!4,colframe=teal!35!black,title={Before you move on}]
What does \mw{भलो} mean? (\textquotedblleft{}Good-natured\textquotedblright{}.) Say it once more.
\end{tcolorbox}
